% SPDX-License-Identifier: CC-BY-NC-ND-4.0
% Copyright (c) 2026 Aymix — workbook content. See LICENSE-CONTENT.
% ============================ APPENDICES ============================
\newcommand{\appheading}[3]{%
  \clearpage\phantomsection
  \addcontentsline{toc}{chapter}{\texorpdfstring{\textbf{Appendix #1}\, — #2}{Appendix #1 — #2}}%
  \markboth{Appendix #1 \ \textbullet\ \ #2}{}%
  \begin{tcolorbox}[enhanced, breakable=false, colback=cBrand, colframe=cBrandDeep,
      boxrule=0pt, arc=2.2mm, left=6mm,right=6mm,top=5mm,bottom=5.5mm,
      overlay={\node[anchor=north east, font=\sffamily\bfseries, text=white!12,
         scale=4.6] at ([xshift=-3mm,yshift=2mm]frame.north east) {#1};}]
    {\sffamily\bfseries\footnotesize\textcolor{cAccent}{\MakeUppercase{Appendix #1}}}\\[1.6mm]
    {\sffamily\bfseries\huge\textcolor{white}{#2}}\\[2.2mm]
    {\sffamily\normalsize\textcolor{white!88}{#3}}
  \end{tcolorbox}\par\vspace{3mm}%
}
\newcommand{\qa}[2]{%
  \par\addvspace{2.2mm}%
  \noindent{\sffamily\bfseries\color{cInterview}\faAngleRight\ #1}\par
  \noindent#2\par}
\newcommand{\gterm}[2]{\par\addvspace{1.4mm}\noindent{\sffamily\bfseries\color{cBrandDeep}#1}\ \textemdash\ #2\par}

% ---------------------------------------------------------------- APPENDIX A
\appheading{A}{DevOps Interview Cheat Sheet}{The mechanism behind the tool --- review this before every interview}

Interviewers at the junior/mid level test \emph{judgment} as much as tool syntax: the ability to explain why a
decision was made, not recite a flag. Every answer below states the mechanism. Read a question, answer it
out loud from memory, then check yourself.

\wbsub{Linux}
\qa{What is PID 1 responsible for?}{PID 1 (\code{systemd} on modern distros) is the first userspace process the kernel starts. It reaps orphaned processes and starts every other service in dependency order.}
\qa{Hard link vs symlink?}{A hard link is a second directory entry pointing at the same inode --- the same file, and it survives the original being deleted. A symlink is a separate file containing a path, and it breaks if the target moves.}
\qa{Why does SIGKILL risk data corruption that SIGTERM does not?}{SIGKILL is handled by the kernel and cannot be caught, so the process gets no chance to flush buffers, close files or release locks before dying.}
\qa{What does load average actually mean?}{The average number of processes running or waiting on CPU/IO over 1/5/15 minutes. A load average above the core count means work is queued.}

\wbsub{Git}
\qa{Merge vs rebase --- the mechanical difference?}{Merge creates a new commit joining two histories, rewriting nothing. Rebase replays commits onto a new base, generating new commit hashes and a linear history.}
\qa{Why is \code{-{}-force-with-lease} safer than \code{-{}-force}?}{It refuses to push if the remote branch has commits you have not fetched, preventing you from silently overwriting a teammate's work.}
\qa{How do you find which commit introduced a bug?}{\code{git bisect} --- a binary search between a known-good and known-bad commit, testing each candidate until it isolates the exact one.}

\wbsub{Docker}
\qa{Image vs container?}{An image is a read-only, layered template; a container is a running (or stopped) instance of that image plus one writable layer on top.}
\qa{Why do multi-stage builds reduce attack surface, not just size?}{Build tools, compilers and dev dependencies never ship in the final image --- fewer packages for a scanner to flag, and fewer tools available to an attacker who gets a shell.}
\qa{CMD vs ENTRYPOINT?}{\code{ENTRYPOINT} defines the fixed executable that always runs; \code{CMD} supplies default arguments that can be overridden at \code{docker run} time.}

\wbsub{Kubernetes}
\qa{What actually happens when you run \code{kubectl apply}?}{\code{kubectl} sends the desired state to the API server, which persists it to etcd. Controllers then reconcile actual cluster state toward that desired state, asynchronously.}
\qa{Why does a Service keep working when Pods are replaced?}{A Service selects Pods dynamically by label, not by a fixed IP list --- as old Pods die and new ones start, the endpoint list updates automatically.}
\qa{Liveness vs readiness, concretely?}{Liveness failing restarts the container. Readiness failing only removes it from load balancing. Conflating the two turns a slow dependency into a restart storm.}
\qa{What is a DaemonSet for?}{Ensuring exactly one copy of a Pod runs on every (or every matching) node --- used for node-level agents such as log shippers and monitoring agents.}

\wbsub{Terraform}
\qa{Why can Terraform not work correctly without state?}{State is the only record mapping configuration blocks to the real IDs of what was created. Without it, Terraform cannot tell what exists from what needs creating.}
\qa{What does \code{terraform plan} actually do?}{Refreshes real resource state, computes the diff against the desired configuration, and prints exactly what would change --- without touching anything.}
\qa{Why remote state with locking instead of local?}{A team needs one source of truth. Without a lock, two simultaneous applies can corrupt the state file or make conflicting changes.}

\wbsub{CI/CD}
\qa{Continuous Delivery vs Continuous Deployment?}{Delivery means every passing build is one click from production. Deployment removes the click and ships automatically.}
\qa{Why tag images by commit SHA instead of \code{latest}?}{Every deployed artifact is traceable to the exact source commit, and rollback means redeploying a known-good tag rather than guessing what \code{latest} currently points at.}
\qa{Why promote the same image across environments instead of rebuilding?}{Rebuilding risks a different dependency resolution or base-image patch sneaking in between builds --- exactly the drift a tested pipeline exists to eliminate.}

\wbsub{AWS \& Cloud}
\qa{Why put a database in a private subnet?}{Private subnets have no route to an internet gateway, so nothing on the internet can reach the database directly regardless of security-group rules. It is a second layer of defence, not the only one.}
\qa{Why prefer IAM roles over access keys for compute identity?}{Roles issue short-lived, automatically rotated credentials scoped to the resource. A static access key is a permanent secret that works from anywhere the moment it leaks.}
\qa{Multi-AZ vs read replica?}{Multi-AZ is synchronous replication for automatic failover (availability). A read replica is asynchronous replication for scaling reads. Different problems --- often combined.}

\wbsub{Monitoring \& Observability}
\qa{Why does p99 latency matter more than average latency?}{Averages smooth outliers into invisibility. p99 exposes the tail of genuinely slow requests, which is where real user pain and SLO breaches live.}
\qa{What is an error budget, practically?}{The allowed amount of SLO violation over a period. Once it is spent, team priority should shift from shipping features to reliability work.}
\qa{Counter vs gauge?}{A counter only increases (resetting on restart); a gauge moves freely up and down. Using a gauge where a counter belongs breaks \code{rate()} calculations.}

\wbsub{Security}
\qa{First move when a secret is committed and pushed?}{Treat it as compromised immediately --- rotate or revoke it --- \emph{then} separately clean git history. Deleting the file in a later commit does not remove it from history.}
\qa{Why run SAST, dependency scanning and image scanning as three separate gates?}{Each catches a different class of risk: insecure code you wrote, known CVEs in libraries you depend on, and known CVEs in the OS packages baked into the image. No single scan covers all three.}
\qa{What does a Kubernetes NetworkPolicy actually restrict?}{Pod-to-pod traffic by label selector. By default every Pod can reach every Pod; a NetworkPolicy is what limits lateral movement if one is compromised.}

\wbsub{System Design \& Troubleshooting}
\qa{How would you scale a read-heavy API?}{Add a cache in front of the database, add read replicas, and consider CDN caching for public/static data before scaling the database vertically.}
\qa{The app is intermittently returning 5xx --- what is your process?}{Check dashboards for error-rate/latency spikes correlated with a deploy or traffic change; check logs by correlation ID; check downstream dependency health; check resource saturation (CPU, memory, connection pool, disk).}
\qa{How do you safely roll back a bad deploy under pressure?}{Redeploy the last known-good, SHA-tagged image. Never rebuild from source under incident pressure --- which is exactly why CI/CD promotes immutable, tagged artifacts.}
\qa{Vertical vs horizontal scaling?}{Vertical is a bigger machine: simple, but with a ceiling and a single point of failure. Horizontal is more machines behind a load balancer, which requires statelessness to work.}

\begin{seniornote}
The candidates who get offers are not the ones who recite the most flags. They are the ones who can say
``here is the trade-off I made, and here is what I would do differently at ten times the scale.'' That sentence
is worth more than a perfect, over-engineered capstone. Practise answering every question above with a
\emph{because}.
\end{seniornote}

% ---------------------------------------------------------------- APPENDIX B
\appheading{B}{The CloudBase Build-Along Master Tracker}{One platform, fourteen increments --- tick each as you finish it}

Your single source of truth for the capstone. By Day 14, \code{terraform apply} should bring up the whole
environment from a clean state and the pipeline should deploy into it unattended.

\wbsub{Phase 1 --- Foundations (Days 1--3)}
\begin{checklist}
  \item \textbf{Day 1:} Base host --- strict-mode \code{bootstrap.sh}, dedicated non-root service user, \code{/opt/cloudbase} with correct modes, systemd unit with a restart policy, verified via \code{journalctl}.
  \item \textbf{Day 2:} Repo conventions --- branch protection on \code{main} (passing CI + one review), \code{CODEOWNERS}, commit convention, pre-commit hook blocking secret patterns.
  \item \textbf{Day 3:} Network design on paper --- one /16 VPC, public and private subnets across two AZs, a bastion as the single audited SSH entry point, least-privilege security groups.
\end{checklist}

\wbsub{Phase 2 --- Containers \& Orchestration (Days 4--6)}
\begin{checklist}
  \item \textbf{Day 4:} Hardened multi-stage Dockerfile --- non-root \code{USER}, pinned base tag, \code{.dockerignore}, layers ordered for cache reuse.
  \item \textbf{Day 5:} Kubernetes deployment --- Deployment with resource requests/limits, Service, ConfigMap and Secret, separated liveness and readiness probes.
  \item \textbf{Day 6:} Production Kubernetes --- Helm chart with per-environment values, Ingress with TLS termination, HPA, least-privilege RBAC.
\end{checklist}

\wbsub{Phase 3 --- Infrastructure as Code \& Delivery (Days 7--10)}
\begin{checklist}
  \item \textbf{Day 7:} Terraform --- VPC, subnets, IGW, NAT gateway and route tables in code, with a remote backend and state locking.
  \item \textbf{Day 8:} Ansible --- an idempotent role replacing \code{bootstrap.sh}, per-environment inventory, secrets in Vault.
  \item \textbf{Day 9:} CI/CD --- test $\rightarrow$ scan $\rightarrow$ build $\rightarrow$ push, SHA-tagged images, secrets from the CI store, manual gate before production.
  \item \textbf{Day 10:} AWS --- ALB as the only public entry, cluster compute, RDS in private subnets, least-privilege IAM roles rather than static keys.
\end{checklist}

\wbsub{Phase 4 --- Operate It Like Production (Days 11--14)}
\begin{checklist}
  \item \textbf{Day 11:} Observability --- Prometheus and Grafana, structured logs with correlation IDs, one real SLO with an alert tied to it.
  \item \textbf{Day 12:} Security --- secrets in a manager or Vault, image scanning as a blocking gate, default-deny NetworkPolicies, IAM/RBAC audit.
  \item \textbf{Day 13:} GitOps --- Argo CD (or Flux) reconciling from git, a canary rollout gated on the Day 11 metrics, and a \emph{rehearsed} rollback.
  \item \textbf{Day 14:} Capstone --- full assembly, a disaster-recovery restore drill, an operational runbook, and an architecture review.
\end{checklist}

\begin{buildalong}[Final completion checklist]
\begin{checklist}
  \item All 13 prior days' concepts are represented somewhere in this infrastructure.
  \item \code{terraform apply} brings up the full environment from a clean state.
  \item The CI/CD pipeline is green and gates on tests plus security scans.
  \item A dashboard and at least one real alert exist and reflect actual SLOs.
  \item Secrets exist in a secrets manager --- never in git, never in the image.
  \item A rollback has been rehearsed, not just designed.
  \item The README explains architecture \emph{decisions} and trade-offs, not just how to run it.
\end{checklist}
\end{buildalong}

% ---------------------------------------------------------------- APPENDIX C
\appheading{C}{Command \& Manifest Quick Reference}{The commands you will actually reach for, grouped by tool}

\wbsub{Linux triage}
\begin{bashbox}[the first five minutes of an incident]
uptime                                  # load average: 1/5/15 min
ps aux --sort=-%cpu | head              # top CPU consumers
df -h && sudo lsof +L1                  # disk full? deleted-but-open files?
journalctl -u <unit> -p err --since "1 hour ago"
ss -tulpn                               # what is listening on which port
\end{bashbox}

\wbsub{Git recovery}
\begin{bashbox}[when something went wrong]
git reflog                              # every HEAD position - your undo history
git revert <sha>                        # safe on shared branches
git bisect start / bad / good <sha>     # binary-search a regression
git push --force-with-lease             # never bare --force
\end{bashbox}

\wbsub{Docker}
\begin{bashbox}[build, inspect, debug]
docker build -t app:$(git rev-parse --short HEAD) .
docker history <image>                  # see layer sizes - find the bloat
docker run --rm -it --entrypoint sh <image>
docker logs -f --tail 100 <container>
docker system df && docker system prune # reclaim disk
\end{bashbox}

\wbsub{Kubernetes}
\begin{bashbox}[the debugging loop]
kubectl get pods -o wide --sort-by=.status.startTime
kubectl describe pod <pod>              # EVENTS at the bottom - read these first
kubectl logs <pod> -c <container> --previous   # logs from the crashed instance
kubectl exec -it <pod> -- sh
kubectl get events --sort-by=.lastTimestamp
kubectl rollout status/undo deployment/<name>
\end{bashbox}

\wbsub{Terraform}
\begin{bashbox}[the safe loop]
terraform init -backend-config=env/prod.hcl
terraform plan -out=tfplan              # ALWAYS plan to a file
terraform apply tfplan                  # apply exactly what you reviewed
terraform state list / show <addr>
terraform force-unlock <lock-id>        # only when you are certain
\end{bashbox}

\wbsub{Ansible}
\begin{bashbox}[idempotency in practice]
ansible-inventory --list -i inventory/prod
ansible-playbook site.yml --check --diff   # dry run, show changes
ansible-playbook site.yml --limit web --tags nginx
ansible-vault edit group_vars/prod/vault.yml
\end{bashbox}

% ---------------------------------------------------------------- APPENDIX D
\appheading{D}{Glossary of Key Terms}{The vocabulary of a DevOps engineer, in plain language}

\begin{multicols}{2}
\small
\gterm{Namespace (Linux)}{Kernel feature giving a process an isolated \emph{view} of a resource (PIDs, network, mounts). One half of what a container is.}
\gterm{cgroup}{Kernel feature limiting and accounting a process group's CPU, memory and IO. The other half of a container.}
\gterm{PID 1}{First userspace process; reaps orphans and supervises services. In a container, your app usually is PID 1.}
\gterm{SIGTERM / SIGKILL}{Catchable, graceful stop versus uncatchable kernel kill. Exit 137 = 128 + 9 = SIGKILL.}
\gterm{Idempotency}{Running an operation again produces the same state. The property that makes automation safe to re-run.}
\gterm{Image layer}{An immutable, content-addressed filesystem diff. Cached and shared between images.}
\gterm{Multi-stage build}{A Dockerfile that builds in one stage and copies only artifacts into a lean runtime stage.}
\gterm{Control plane}{Kubernetes API server, scheduler, controller manager and etcd --- the brain that records and reconciles desired state.}
\gterm{Reconciliation loop}{A controller continuously driving actual state toward desired state. The core idea of Kubernetes and GitOps.}
\gterm{Liveness vs readiness}{Liveness failure restarts the container; readiness failure only removes it from load balancing.}
\gterm{Service / Ingress}{A Service gives stable in-cluster networking by label selector; an Ingress is an HTTP reverse proxy for external traffic, usually terminating TLS.}
\gterm{HPA}{Horizontal Pod Autoscaler --- scales replica count from a metric ratio.}
\gterm{Helm chart}{A templated, versioned package of Kubernetes manifests with per-environment values.}
\gterm{Terraform state}{The map from configuration to real resource IDs. Contains secrets --- never commit it.}
\gterm{Plan / apply}{Plan computes and prints the diff without changing anything; apply executes exactly that diff.}
\gterm{Drift}{Reality diverging from what the code declares, usually because someone changed it by hand.}
\gterm{Remote backend}{Shared, locked state storage so a team cannot corrupt state with concurrent applies.}
\gterm{CIDR}{\code{10.0.0.0/16} --- the prefix length fixes the network portion; the rest addresses hosts. /24 = 254 usable.}
\gterm{Security group vs NACL}{Stateful, instance-level, allow-only versus stateless, subnet-level, ordered with explicit deny.}
\gterm{NAT gateway}{Lets private-subnet hosts reach the internet outbound without being reachable inbound. Bills hourly plus per GB.}
\gterm{Bastion / jump host}{A single hardened, audited SSH entry point to hosts with no public route.}
\gterm{IAM role vs access key}{A role issues short-lived rotated credentials to a workload; a static key is a permanent secret that works from anywhere.}
\gterm{Least privilege}{Grant exactly the permissions needed, no more. Explicit deny always wins in IAM evaluation.}
\gterm{TLS termination}{Decrypting HTTPS at the load balancer or ingress so backends handle plain HTTP inside the trust boundary.}
\gterm{DNS TTL}{How long resolvers cache an answer. Lower it days before a planned cutover, never during one.}
\gterm{Artifact promotion}{Building once and moving the same immutable, SHA-tagged image through environments.}
\gterm{GitOps}{Git as the single source of truth, with an agent continuously reconciling the cluster to match it.}
\gterm{Blue-green / canary}{Two full environments with an instant switch, versus shifting a small traffic percentage while watching metrics.}
\gterm{SLI / SLO / error budget}{The measured indicator, the target for it, and the amount of failure you may spend before reliability work takes priority.}
\gterm{p99 latency}{The 99th percentile --- the slow tail where real user pain lives, invisible in an average.}
\gterm{Counter vs gauge}{Monotonically increasing versus free-moving. Using the wrong one breaks \code{rate()}.}
\gterm{Cardinality}{Number of distinct label combinations in a metric. High cardinality is the usual cause of a melted Prometheus.}
\gterm{NetworkPolicy}{Label-selector firewall for pod-to-pod traffic; without one, everything can reach everything.}
\gterm{RTO / RPO}{How fast you must recover, and how much data you may lose. Both are business decisions, not technical ones.}
\end{multicols}

\vspace{4mm}
\begin{center}\small\itshape\color{cInk!65}
Good luck --- now go build something, and break it on purpose before production does it for you.
\end{center}
